\documentclass[preprint,12pt]{elsarticle}

\begin{document}

\begin{frontmatter}

\title{Claim-Safe Artifact Engineering for Software-Only EDA Experiments: Design and Evaluation of ArtifactGate-EDA}

\begin{abstract}
ArtifactGate-EDA is a software engineering method and tool for preparing,
validating, replaying, and auditing software-only EDA experiment artifacts. The
tool indexes artifacts through adapters, records file hashes and provenance,
assigns bounded evidence levels, checks claims against a policy, builds replay
manifests, detects corrupted packages, and emits reviewer-facing reports. The
IST evaluation covers multi-adapter ingestion, replay reproducibility, negative
claim injection, corrupted artifact detection, evidence-level classification,
external case generalization, scalability, baseline comparison, ablation, and
optional-backend audit experiments. The package also includes a generated
reviewer-walkthrough dry run for author-side protocol preparation; this is not
used as human or expert-study evidence.
ArtifactGate-EDA is a software engineering method and tool for making
software-only EDA experiment artifacts claim-safe, replay-checkable, and
reviewer-auditable.
\end{abstract}

\begin{keyword}
artifact evaluation \sep reproducibility \sep electronic design automation \sep software engineering \sep claim checking
\end{keyword}

\end{frontmatter}

\section{Introduction}

EDA artifact packages need more than file availability. Reviewers need to know
which claims are supported by the included software artifacts and which claims
would require stronger evidence outside the package. ArtifactGate-EDA addresses
this as a claim-safe artifact engineering problem for software-only EDA
experiments.

\section{Related Work}

The final submission should add author-verified references for artifact
evaluation, reproducible software packaging, EDA experiment reproducibility,
claim-evidence traceability, and empirical software engineering evaluation.

\section{Method}

ArtifactGate-EDA represents each file as an artifact record containing a
relative path, artifact type, SHA-256 hash, tool label, adapter name, evidence
level, claim binding, unsupported boundary, and creation command. The method
combines adapter ingestion, validation, replay packaging, claim checking, and
report generation.

\section{Tool Implementation}

The command-line interface exposes \texttt{ingest}, \texttt{validate},
\texttt{replay}, \texttt{claim-check}, \texttt{report}, \texttt{package},
\texttt{compare}, \texttt{benchmark}, and \texttt{ablate}. Core reproduction
does not require proprietary local tools.

\section{Evaluation Design}

The evaluation is organized as RQ1--RQ10: multi-adapter ingestion, replay
reproducibility, ClaimBench-EDA claim checking, corrupted artifact detection,
evidence-level classification, external case generalization, scalability,
baseline comparison, ablation, optional-backend audit, and generated reviewer
dry-run preparation. The full command is \texttt{make ist-all}. The
ClaimBench-EDA gold standard is policy-seeded rather than independently
human-rated. The replay repeatability extension records 10 repeated runs for
each core replay case and reports local Docker availability separately.

\section{Results}

The generated reports show that all current RQ report files are present. RQ1 is
summarized in \texttt{reports/rq1\_multi\_adapter\_ingestion.md}. RQ2 is
summarized in \texttt{reports/rq2\_replay\_reproducibility.md} and
\texttt{reports/rq2\_replay\_summary.csv}; the repeatability extension is
summarized in \texttt{reports/rq2\_replay\_repeats.md},
\texttt{reports/rq2\_replay\_repeat\_summary.csv}, and
\texttt{reports/rq2\_docker\_ci\_status.csv}. The latter records Docker
availability without claiming a local Docker pass. RQ3 evaluates 1300 ClaimBench-EDA
rows in \texttt{reports/rq3\_claimbench\_summary.md},
\texttt{reports/rq3\_confusion\_by\_claim\_type.csv}, and
\texttt{reports/rq3\_claimbench\_results.csv}, with zero critical false
negatives in the current generated summary. RQ4 in
\texttt{reports/rq4\_corrupted\_artifact\_detection.md} evaluates a
deterministic software-only corruption benchmark with 900 corrupted-artifact
instances across 30 defect classes, plus 30 clean specificity cases, with the
confusion matrix in
\texttt{reports/rq4\_error\_classification\_matrix.csv}. The extended
taxonomy, instance results, clean specificity set, and severity-weighted score
are reported in \texttt{reports/rq4\_defect\_taxonomy.md},
\texttt{reports/rq4\_corruption\_extended\_results.csv},
\texttt{reports/rq4\_clean\_specificity\_cases.csv}, and
\texttt{reports/rq4\_severity\_weighted\_detection.md}. RQ5 classifies 1000
evidence-level rows with holdout and cross-tool holdout labels in
\texttt{reports/rq5\_evidence\_level\_classification.md} and
\texttt{reports/rq5\_evidence\_level\_confusion\_matrix.csv}, with split
metrics in \texttt{reports/rq5\_evidence\_level\_holdout\_summary.csv}. E6/G9
reports 10 public external cases, 8 open-source ArtifactGate
replay-package/validation cases, adapter success rate 1.00, observed manual
adapter fix rate 0, and observed hardware dependency count 0
in \texttt{reports/rq6\_external\_case\_generalization.md},
\texttt{reports/rq6\_external\_case\_generalization.csv}, and
\texttt{reports/rq6\_external\_case\_generalization\_summary.csv}. The
scalability extension writes and row-parses synthetic ArtifactGate manifests up
to 100k rows, with 53 measured manifest-processing runtime observations across
seven scales and linear-fit R2 above 0.99 in
\texttt{reports/rq7\_scalability\_summary\_extended.md},
\texttt{reports/rq7\_scalability\_extended\_runtime.csv},
\texttt{reports/rq7\_scalability\_extended\_memory.csv}, and
\texttt{reports/rq7\_scalability\_model\_fit.csv}. The baseline extension runs
seven deterministic baseline/target emulators on one shared software-only
fixture. Each emulator writes method-specific outputs under
\texttt{outputs/rq7\_baseline/}, and the evaluator reads those generated outputs
to produce 56 task observations in
\texttt{reports/rq7\_baseline\_task\_results.csv} and summary metrics in
\texttt{reports/rq7\_baseline\_task\_execution.md}. RQ8 generates ten
ablation variants with observation-level rows, bootstrap confidence intervals,
and effect sizes in \texttt{reports/rq8\_ablation.md},
\texttt{reports/rq8\_ablation\_effect\_sizes.csv}, and
\texttt{reports/rq8\_ablation\_observations.csv}. RQ7, RQ8, and RQ9 are
anchored in \texttt{reports/rq7\_baseline\_comparison.md},
\texttt{reports/rq8\_ablation.md}, and
\texttt{reports/rq9\_local\_backend\_audit.md}. RQ10 generates a
reviewer-walkthrough dry run with 16 task observations and command-log
timestamps in \texttt{reports/rq10\_reviewer\_walkthrough.md},
\texttt{reports/rq10\_reviewer\_walkthrough\_observations.csv}, and
\texttt{reports/rq10\_reviewer\_walkthrough\_command\_log.csv}. It is
preparatory material for an author-side expert walkthrough, not completed
expert-walkthrough evidence, a multi-participant human-subject study, or
measured human timing result.

\section{Discussion}

The evaluation supports the software engineering claim that ArtifactGate-EDA
adds a claim-boundary layer beyond ordinary file bundling, checksums, shell
scripts, and simple manifests.

\section{Threats to Validity}

The claim dataset and corrupted-artifact suite are controlled and deterministic.
Future work should add independent manuscripts, additional artifact formats, and
externally reviewed defect cases. The external-case generalization benchmark
uses public source files and ArtifactGate replay-packaging/validation or
metadata-only ingestion checks; it does not claim independent EDA tool
execution, circuit correctness, hardware validation, or vendor timing closure.
The scalability memory values are estimates derived from manifest size and row
count, not measured process RSS. The baseline benchmark uses a deterministic
task harness and operator-step time estimates, not a human-subject timing
study. The generated reviewer-walkthrough dry run is not evidence that the
optional G13 expert walkthrough has been completed; participant or
expert-walkthrough claims require separate author-confirmed human data.

\section{Limitations}

ArtifactGate-EDA checks artifact-level support for bounded claims. It does not
prove design correctness and does not convert software simulation or synthesis
artifacts into stronger device-side evidence. The current evidence ceiling is
\texttt{L4\_SYNTHESIS}.

\section{Reproducibility and Data Availability}

The SoftwareX baseline is archived with DOI \texttt{10.5281/zenodo.20789516}.
The IST evaluation snapshot is archived as version \texttt{v0.1.3} with DOI
\texttt{10.5281/zenodo.20798200}. The IST evaluation layer is regenerated with
\texttt{make install} and \texttt{make ist-all}. The IST archive is a software
package and generated-data snapshot. It does not add hardware validation,
vendor timing closure, bitstream, DFX deployment, or board-validation evidence.

\end{document}
